% ----------------------------------------------------------------------------
%                           FUNDAMENTAÇÃO TEÓRICA
% ----------------------------------------------------------------------------

\chapter{CONCEITOS GERAIS E REVISÃO DA LITERATURA}
\label{chap:fundamentacao}

Neste capítulo, será apresentada uma breve introdução para à Doença de Parkinson, abordando como um sistema de dados pode auxiliar no apoio à decisão na área da saúde e a importância do teste de usabilidade do sistema.

\section{Doença de Parkinson}
\label{sec:Parkinson}

A Doença de Parkinson (DP) é uma afecção neurodegenerativa crônica e progressiva do sistema nervoso central, configurando-se como a segunda doença neurodegenerativa mais comum associada ao envelhecimento \cite{magalhaes2022teorias}. Independente da idade de início, a fisiopatologia primária dos sintomas decorre da degeneração dos neurônios dopaminérgicos localizados na parte compacta da substância negra mesencefálica \cite{silva2021doenca, Santos2024, Lima2023}. Estima-se que, ao manifestarem-se os primeiros sinais motores, cerca de 60\% desses neurônios já tenham sido previamente degradados \cite{silva2021doenca}.

A redução substancial e continuada da dopamina, neurotransmissor essencial para a regulação dos circuitos dos núcleos da base e modulação do córtex motor, desencadeia o surgimento de uma sintomatologia complexa e heterogênea \cite{Lima2023, Cabreira2019}. Clinicamente, a DP manifesta-se por meio de uma tétrade motora cardinal, composta por bradicinesia, rigidez muscular, tremor de repouso e alterações posturais e de marcha \cite{Santos2024, Luiz2022}. Além dos sintomas motores, a literatura oficial reconhece que a apresentação clínica engloba disfunções não motoras severas de ordem cognitiva, autonômica e comportamental, compreendendo alterações do olfato, distúrbios do sono, constipação, sintomas psicóticos e evolução para demência \cite{brasil2025, Santos2024, Cabreira2019}.

Existe uma prevalência de 0,3\% na população geral, sendo mais comumente na população acima dos 80 anos. Em termos epidemiológicos, a taxa de incidência global de DP em pessoas com 40 anos ou mais apresenta diferenças significativas entre os sexos, sendo de 37,55 por 100.000 pessoas-ano em mulheres e 61,21 por 100.000 pessoas-ano em homens. As taxas de incidência aumentam progressivamente com a idade; entre as mulheres, variam de 3,26 por 100.000 pessoas-ano na faixa etária de 40 a 49 anos, para 103,48 com mais de 80 anos. Entre os homens, esse aumento é ainda mais expressivo, saltando de 3,57 por 100.000 pessoas-ano (na faixa de 40 a 49 anos) para 258,47 na idade superior a 80 anos. Uma variedade clínica da doença acomete também pacientes antes dos 40 anos.

Não existe cura para a doença, sendo o tratamento estritamente voltado para o controle de sintomas. O manejo envolve abordagens não farmacológicas, fisioterapia, fonoterapia e terapia ocupacional, aliados ao tratamento medicamentoso. Embora a progressão da doença seja individual e possua um curso clínico variável, o manejo adequado exige a estratificação dos pacientes em três subtipos principais para fins de prognóstico: (1) indivíduos com predominância motora leve, que possuem progressão lenta e boa resposta medicamentosa; (2) indivíduos com subtipo intermediário; e (3) indivíduos com subtipo maligno difuso, caracterizados por sintomas precoces proeminentes, progressão rápida e má resposta terapêutica \cite{brasil2025}.

Em virtude do caráter evolutivo da enfermidade e da ausência de biomarcadores laboratoriais ou exames de imagem definitivos para a confirmação diagnóstica \textit{in vivo}, a conduta médica permanece estritamente dependente da avaliação semiológica clínica contínua \cite{Lima2023, Cabreira2019}.


\subsection{Estadiamento clínico e evolução funcional}
\label{subsec:estadiamento-clinico}

Para caracterizar as fases de comprometimento motor, a escala de Hoehn e Yahr (HY) classifica a condição funcional observada em cinco estágios. Esses estágios descrevem a gravidade clínica no momento da avaliação; não representam intervalos de duração fixa nem uma sequência obrigatória para todos os pacientes \cite{Goetz2004HY}.

\begin{enumerate}
    \item \textbf{Estágio 1:} comprometimento unilateral, geralmente com limitação funcional mínima ou ausente.
    \item \textbf{Estágio 2:} comprometimento bilateral ou axial, sem alteração do equilíbrio.
    \item \textbf{Estágio 3:} comprometimento bilateral com instabilidade postural, mantendo independência física.
    \item \textbf{Estágio 4:} incapacidade acentuada, embora ainda seja possível permanecer em pé ou caminhar sem auxílio.
    \item \textbf{Estágio 5:} dependência de cadeira de rodas ou permanência no leito, salvo quando há assistência.
\end{enumerate}

Essa síntese baseia-se na revisão da MDS \cite{Goetz2004HY}. A HY enfatiza lateralidade, equilíbrio e independência, sem abranger adequadamente manifestações não motoras. Sua natureza ordinal impede interpretar distâncias entre estágios como equivalentes. Assim, o estadiamento deve ser complementado por uma avaliação multidimensional, como a MDS-UPDRS.

\subsection{A Escala MDS-UPDRS e a Avaliação Clínica} 
\label{subsec:mds-updrs}

\paragraph{A Necessidade de Padronização Semiológica} 
A mensuração empírica dos sintomas da Doença de Parkinson, quando desprovida de critérios estandardizados, introduz uma elevada subjetividade e variabilidade inter e intraexaminadores na estimativa do comprometimento motor \cite{Luiz2022}. Torna-se, portanto, de vital necessidade a adoção de métricas clínicas padronizadas e validadas internacionalmente \cite{Goetz2008}, assegurando que equipes multiprofissionais e diferentes especialistas avaliem a progressão motora, a severidade dos sintomas e a eficácia das intervenções terapêuticas de maneira coesa, objetiva e reprodutível \cite{Santos2024, Luiz2022}. 

\paragraph{A Estrutura Metodológica da Escala MDS-UPDRS} 
Para suprir essa demanda por padronização clinimétrica, a \textit{Movement Disorder Society} (MDS) promoveu uma revisão estruturada da tradicional \textit{Unified Parkinson's Disease Rating Scale} (UPDRS), consolidando a escala MDS-UPDRS como o padrão-ouro mundial para a quantificação da gravidade e progressão da DP \cite{Goetz2008, Luiz2022}. A escala possui uma organização sistemática dividida em quatro domínios de avaliação complementares \cite{Santos2024, Goetz2008}: 

\begin{itemize} 
    \item \textbf{Parte I (Experiências Não Motoras da Vida Diária):} Avalia o impacto de aspectos não motores na rotina do paciente, tais como comprometimento cognitivo, depressão, ansiedade, distúrbios do sono e alucinações \cite{Santos2024, Goetz2008}; 
    \item \textbf{Parte II (Experiências Motoras da Vida Diária):} Mensura a percepção do paciente ou de seu cuidador quanto às limitações funcionais provocadas pelos sintomas motores em atividades autônomas, como fala, escrita, vestuário, mastigação e higiene \cite{Santos2024, Goetz2008}; 
    \item \textbf{Parte III (Exame Motor):} Consiste na avaliação clínica objetiva realizada pelo especialista, destinada a quantificar a intensidade das manifestações motoras ativas do paciente \cite{Santos2024, Goetz2008, Luiz2022}; 
    \item \textbf{Parte IV (Complicações Motoras):} Examina o surgimento de complicações secundárias induzidas pela farmacoterapia dopaminérgica de longo prazo, abrangendo discinesias involuntárias e flutuações motoras, caracterizadas pela alternância entre períodos de melhor controle dos sintomas sob efeito da medicação (\textit{ON}) e períodos de redução desse benefício (\textit{OFF}) \cite{Santos2024, Goetz2008}.
\end{itemize} 

Dentre estes domínios, a \textbf{Parte III} configura-se como o núcleo fundamental de quantificação do estado motor do indivíduo \cite{Lima2023, Luiz2022}. Sua aplicação envolve um protocolo de tarefas padronizadas para mensurar detalhadamente os sinais cardinais da doença. A bradicinesia, sinal obrigatório para a caracterização do parkinsonismo \cite{Lima2023, Cabreira2019}, é rigorosamente avaliada por meio de testes cinemáticos de membros superiores e inferiores, como a tarefa de pinça ou batida dos dedos (\textit{finger tapping test}), abertura e fechamento das mãos, movimentos de pronação-supinação do antebraço e agilidade das pernas \cite{Lima2023, Luiz2022}. Nesses exames, o clínico analisa a velocidade, a amplitude, o ritmo e o declínio progressivo do movimento no tempo \cite{Lima2023, Luiz2022}. Simultaneamente, a Parte III afere a rigidez muscular por meio da movimentação passiva das articulações e mensura os tremores de repouso e posturais/cinéticos \cite{Santos2024, Goetz2008}. Cada item da escala é pontuado segundo um escore ordinal de 5 pontos ($0 = \text{normal}$, $1 = \text{discreto}$, $2 = \text{leve}$, $3 = \text{moderado}$ e $4 = \text{grave}$) \cite{Santos2024, Goetz2008, Luiz2022}. 

\paragraph{Fundamentação clinimétrica e interpretação dos escores}
O estudo de validação da MDS-UPDRS incluiu 877 pacientes de 39 centros e examinou consistência interna e estrutura fatorial. Os resultados sustentaram a interpretação de cada parte separadamente, sem recomendar a soma das quatro partes como um único desfecho clínico \cite{Goetz2008}. Portanto, os painéis devem apresentar os escores por domínio, preservando a identificação da parte avaliada.

O exame motor também exige contextualização: o formulário registra o estado \textit{ON/OFF} associado à avaliação da Parte III \cite{Goetz2008}. Comparações longitudinais devem considerar essa informação, pois avaliações realizadas em condições terapêuticas distintas não são diretamente intercambiáveis. A pontuação representa o comprometimento observado nas condições de coleta; sua variação, isoladamente, não demonstra progressão neurodegenerativa nem eficácia de uma intervenção.

Como requisito de projeto derivado dessa fundamentação, propõe-se armazenar data da avaliação, versão do instrumento, estado medicamentoso, respostas individuais e escores por parte. Campos não preenchidos devem permanecer identificados como ausentes, evitando sua conversão automática em zero. Esses requisitos orientam a rastreabilidade dos dados e precisam ser verificados na implementação e na validação do sistema.

\paragraph{A Ligação Tecnológica e Computacional: O Painel Parkinson Desktop} 
A despeito da robustez clinimétrica da MDS-UPDRS, sua aplicação tradicional baseada em formulários analógicos em papel apresenta importantes limitações operacionais. A coleta manual sujeita o processo a equívocos de transcrição e a erros de somatório aritmético na consolidação dos escores totais de cada parte, o que pode comprometer a precisão do acompanhamento e induzir a interpretações equivocadas do estado do paciente. 

A superação desses gargalos operacionais viabiliza-se pela transição para um ecossistema digital estruturado, consubstanciado na arquitetura proposta para o \textit{Painel Parkinson Desktop}. Desenvolvido para atuar como um Sistema de Apoio à Decisão Clínica (SADC) de alto desempenho, o sistema conecta os requisitos da literatura médica à Ciência da Computação. Ao integrar a interface reativa em VueJS (encapsulada pelo \textit{framework} Electron.js) a um \textit{back-end} local em Python (Flask) e ao banco de dados relacional PostgreSQL (vinculado ao modelo de dados do Sistema Integrado de Dados Biomédicos, SIDABI) \cite{folador2021}, a plataforma automatiza o cálculo instantâneo dos escores da MDS-UPDRS, eliminando por completo os erros de somatório manual e a perda de estado de seleções médicas. 

Sob a perspectiva da engenharia de software e da infraestrutura de dados, a arquitetura local (\textit{desktop}) descentraliza o processamento biométrico, eliminando a latência de rede decorrente de servidores em nuvem e assegurando operação plena em modo \textit{offline}. Essa descentralização garante a persistência segura dos dados locais no ambiente institucional e atende aos requisitos da Lei Geral de Proteção de Dados (LGPD, Lei nº 13.709/2018). Por meio de mecanismos automatizados de anonimização em segundo plano (em conformidade com o Art. 13 da LGPD), os dados clínicos sensíveis permanecem sob controle estrito do ambiente hospitalar, enquanto \textit{hashes} criptográficos viabilizam a sincronização segura com bases centrais de pesquisa sem expor a identidade dos pacientes. Dessa forma, a digitalização local e estruturada possibilita a geração automatizada de históricos longitudinais precisos e gráficos de evolução temporal, fornecendo ao especialista uma ferramenta robusta, ágil e segura para o suporte à decisão clínica sem a dependência de infraestrutura em nuvem.

% !TeX root = ../../../principal.tex
\subsection{Relação entre os dados da MDS-UPDRS e a representação gráfica}
\label{subsec:dados-mds-graficos}

A estrutura multidimensional da MDS-UPDRS orienta a organização dos painéis: cada avaliação associa uma data aos escores das partes I, II, III e IV. Essa associação permite representar tanto o acompanhamento temporal quanto a distribuição dos valores entre domínios clínicos. A recomendação de interpretar as partes separadamente deve ser preservada na visualização, evitando tratar a soma dos domínios como um desfecho clínico validado \cite{Goetz2008}.

\paragraph{Mapeamento dos dados para as coordenadas do gráfico}
No sistema analisado, o histórico do paciente fornece as datas das avaliações e suas respectivas pontuações. Esses dados são associados às categorias do eixo horizontal e aos valores do eixo vertical da série denominada ``Pontuação''. Assim, a data determina a posição de cada avaliação na sequência apresentada, enquanto a pontuação determina a altura do ponto no gráfico \cite{parkinsonDashboardFront2026}.

Assim, cada ponto representa uma associação $(d_k,T_k)$, em que $d_k$ identifica a data da avaliação e $T_k$ é a pontuação total fornecida ao painel. A visualização utiliza esse valor recebido, cuja fórmula de agregação precisa ser verificada na origem dos dados. Além disso, as datas são utilizadas como categorias na ordem recebida, sem ordenação cronológica explícita no tratamento analisado. A distância horizontal entre pontos não deve ser interpretada automaticamente como proporcional ao tempo transcorrido.

\paragraph{Particionamento por domínio e seleção da avaliação}
O particionamento visual corresponde às quatro partes da escala. No histórico, a seleção de um ponto recupera os escores das partes I, II, III e IV associados à avaliação escolhida. Esses valores atualizam quatro barras, uma para cada domínio, acompanhadas da data correspondente \cite{parkinsonDashboardFront2026}. A visualização combina, dessa forma, uma visão temporal com o detalhamento da avaliação selecionada.

Na atualização por seleção, o comprimento preenchido da barra da parte $j$ corresponde a uma transformação percentual:
\begin{equation}
    P_{j,k}=\left\lceil 100\frac{S_{j,k}}{M_j}\right\rceil,
\end{equation}
em que $S_{j,k}$ representa o escore recebido para a parte $j$ na avaliação $k$, $M_j$ é o denominador adotado na visualização e o arredondamento é feito para cima. Os denominadores presentes no sistema analisado são $(52,52,128,24)$. Por exemplo, o valor 26 na Parte I produz uma barra de 50\%. Esse percentual descreve a proporção em relação ao denominador adotado pela interface; não equivale a uma porcentagem de progressão da doença, incapacidade ou resposta terapêutica.

\paragraph{Limitações e requisitos de fidelidade da visualização}
A descrição da implementação deve ser distinguida de sua validação clínica. O denominador 128 da Parte III precisa ser confrontado com a versão do instrumento e os itens efetivamente coletados, antes de ser apresentado como limite da escala. A inicialização das barras utiliza os valores brutos, enquanto a seleção aplica a transformação percentual; essa diferença requer uniformização. O gráfico temporal também fixa o eixo vertical entre 0 e 200, uma configuração da interface que não define a faixa clínica do instrumento \cite{parkinsonDashboardFront2026}.

Outra limitação identificada no tratamento dos dados é a exclusão de valores iguais a zero e a aplicação de critérios de filtragem diferentes ao histórico e aos escores por parte. Isso pode excluir avaliações válidas e comprometer a correspondência entre o ponto selecionado e seus domínios. Como requisito para a evolução do sistema, cada avaliação deve manter um identificador comum entre data, escores e barras, com validação explícita de valores ausentes e preservação dos zeros. Recomenda-se ainda apresentar o acompanhamento por parte, identificar o estado ON/OFF da Parte III e documentar os denominadores utilizados, mantendo a interpretação clínica sob responsabilidade do profissional.


\section{Sistemas de apoio à decisão na saúde}
\label{sec:Sistemas de apoio}

Os sistemas de apoio à decisão clínica são ferramentas computacionais que auxiliam os profissionais da saúde na análise de dados e na tomada de decisões, contribuindo para o processamento de informações complexas no contexto clínico \cite{shortliffe2014}.

Atualmente, os Sistemas de Apoio à Decisão Clínica (SADC) representam abordagens tecnológicas fundamentais para otimizar a prática profissional, pois fornecem instruções que complementam o conhecimento e as habilidades do corpo clínico sem substituí-los. A adoção dessas ferramentas atua diretamente na minimização de erros médicos e na redução de custos e de tempo relacionados ao diagnóstico. Além de funcionarem como um repositório unificado do histórico do paciente, os SADC cruzam grandes volumes de informações para gerar recomendações e alertas personalizados, demonstrando particular eficiência nas fases de prevenção e apoio diagnóstico. Dessa forma, a informatização dos dados aumenta a segurança no manejo clínico e a perspectiva de assertividade nos resultados terapêuticos \cite{silva2022sistemas}.

% !TeX root = ../../../principal.tex
\section{Interoperabilidade de dados clínicos com HL7 FHIR}
\label{sec:hl7-fhir}

A interoperabilidade em saúde envolve a troca de informações entre sistemas e a preservação de seu significado clínico. Um hospital, um laboratório e uma clínica podem utilizar sistemas diferentes, com formas próprias de organizar os dados do paciente. O FHIR estabelece estruturas comuns para que essas aplicações possam compartilhar informações de maneira padronizada, ou seja, ``falar a mesma língua''. No acompanhamento da Doença de Parkinson, isso permite organizar o intercâmbio de avaliações e resultados entre serviços de saúde. O artigo da \citeonline{ctcFHIR2024} contextualiza essa necessidade; a especificação oficial da Health Level Seven International (HL7) fundamenta sua realização técnica.

O \textit{Fast Healthcare Interoperability Resources} (FHIR) é um padrão da HL7 para a troca eletrônica de informações em saúde. Sua unidade básica é o recurso: uma estrutura que representa uma entidade ou informação clínica, com elementos e tipos de dados definidos. Recursos podem ser relacionados por referências, permitindo compor registros que associem o paciente, as avaliações e seus resultados. A especificação contempla representações como JSON e XML, facilitando o processamento das informações por aplicações distintas \cite{hl7FHIRoverview2023}.

\subsection{Troca de informações e preservação do significado}

Uma das formas de intercâmbio previstas pelo FHIR é a interface REST sobre HTTP, com operações para leitura, pesquisa, criação e atualização de recursos. Entretanto, disponibilizar uma API ou transmitir arquivos JSON não é suficiente para caracterizar uma integração FHIR: o conteúdo e as interações precisam respeitar as regras aplicáveis da especificação \cite{hl7FHIRhttp2023}.

Os perfis permitem adaptar os recursos a um contexto de uso, estabelecendo restrições sobre elementos, cardinalidades, terminologias e extensões. Essas definições ajudam os sistemas participantes a concordar sobre quais informações devem existir e como serão interpretadas. Um guia de implementação organiza regras e artefatos para um cenário de integração. Assim, a interoperabilidade exige um acordo sobre versão do FHIR, perfis adotados e códigos utilizados, além da compatibilidade do transporte \cite{hl7FHIRprofiling2023}.

\subsection{Exportação dos dados do paciente em PDF e FHIR}

Para o Painel Parkinson Desktop, a proposta é oferecer ao usuário a escolha do formato de exportação dos dados do paciente selecionado. O PDF apresenta um relatório destinado à leitura, impressão e consulta pelos profissionais, enquanto a exportação FHIR fornece dados estruturados para processamento por outros sistemas compatíveis. As duas opções atendem a finalidades complementares: comunicação por meio de um documento legível e intercâmbio de informações que possam ser incorporadas a outro sistema.

A exportação FHIR pode ser organizada em um arquivo JSON contendo recursos relacionados. O recurso \textit{Patient} representa os dados administrativos e de identificação do paciente, enquanto um \textit{Bundle} permite reunir uma coleção de recursos em uma única estrutura \cite{hl7FHIRpatient2023,hl7FHIRbundle2023}. Na proposta de exportação, essa coleção poderia reunir o paciente e suas avaliações, preservando as referências entre os registros. A importação depende de o sistema destinatário aceitar a versão e os perfis utilizados; o arquivo não garante compatibilidade automática com qualquer aplicação.

\subsection{Representação das avaliações MDS-UPDRS na exportação}

No contexto da exportação proposta, o recurso \textit{Questionnaire} permite definir perguntas, agrupamentos e opções de resposta, enquanto \textit{QuestionnaireResponse} registra as respostas de uma aplicação do questionário. Essa separação possibilita distinguir a definição do instrumento do registro individual de sua aplicação \cite{hl7FHIRquestionnaire2023}.

Como proposta de modelagem para o Painel Parkinson, os grupos do questionário poderiam corresponder às quatro partes da MDS-UPDRS, preservando a identificação dos itens e a versão do instrumento. Essa correspondência é uma proposta para o domínio da aplicação, não um perfil MDS-UPDRS validado ou uma integração já implementada. Sua adoção dependeria do mapeamento dos dados existentes e da definição de regras compartilhadas com os sistemas destinatários.

O recurso \textit{Observation} representa medições e avaliações clínicas e pode registrar resultados de instrumentos de avaliação. Nesse modelo proposto, os escores de cada parte poderiam ser representados por observações separadas, associadas ao paciente, à data da avaliação e ao instrumento de origem. A especificação prevê elementos para identificar o tipo de observação, seu valor, momento clinicamente relevante e referências aos dados dos quais o resultado foi derivado \cite{hl7FHIRobservation2023}.

Para o acompanhamento gráfico, esse mapeamento permitiria recuperar a data e o escore de cada domínio sem depender dos nomes internos de campos de um único sistema. A informação sobre o estado ON/OFF precisaria de representação acordada no perfil adotado, preservando o contexto do exame motor. Também seria necessário manter a distinção entre respostas ausentes e escores iguais a zero, evitando transformar ausência de informação em um resultado clínico normal.

\subsection{Segurança e limites da aplicação do padrão}

A especificação FHIR não é, por si só, um protocolo de segurança. A implantação requer mecanismos de autenticação, autorização, proteção do transporte e auditoria, definidos conforme o ambiente de uso \cite{hl7FHIRsecurity2023}. Portanto, adotar estruturas FHIR não demonstra automaticamente a proteção dos dados ou a conformidade jurídica da aplicação.

Para o Painel Parkinson Desktop, essa fundamentação orienta a proposta de disponibilizar exportação em PDF e FHIR. A concretização da opção FHIR exige selecionar a versão e os perfis de integração, validar os recursos gerados e testar sua importação em sistemas compatíveis. O padrão organiza a representação e o intercâmbio dos dados; a interpretação dos escores e sua utilização no acompanhamento permanecem vinculadas à fundamentação clínica da MDS-UPDRS.


\section{Ciência de dados}
\label{sec:Ciencia de dados}

A Ciência de Dados é uma área multifacetada que não possui uma definição única, podendo ser descrita simultaneamente como uma ciência, uma disciplina, um fluxo de trabalho e uma profissão \cite{mendes2025fundamentos}. Historicamente, ela possui raízes profundas na Estatística, especialmente na Análise Exploratória de Dados, evoluiu para lidar com a análise de grandes e complexos conjuntos de informações, frequentemente chamados de Megadados ou \textit{Big Data} \cite{morettin2019introducao}.  Como disciplina, trata-se de um campo interdisciplinar que integra conhecimentos de Matemática, Estatística, Ciência da Computação e o domínio específico da aplicação estudada \cite{mendes2025fundamentos}. A literatura divide a Ciência de Dados em três perspectivas principais: a estatística, que orienta a coleta e análise dos dados; a computacional, focada no desenvolvimento de algoritmos e no armazenamento em larga escala; e a humana, que busca aplicar esses métodos para resolver problemas reais e comunicar os resultados de forma clara \cite{morettin2019introducao}.  Encarada como um fluxo de trabalho, a Ciência de Dados é um processo contínuo e iterativo que engloba a coleta, limpeza, organização, análise, interpretação e visualização das informações \cite{mendes2025fundamentos}. É importante destacar que ela é mais ampla do que a simples mineração de dados enquanto a mineração foca em descobrir padrões em dados estruturados, a Ciência de Dados engloba todo o ciclo de vida da informação, incluindo dados não estruturados e o uso de aprendizado de máquina \cite{mendes2025fundamentos}. Por fim, devido à sua complexidade, a resolução de problemas nessa área exige o trabalho de equipes multidisciplinares, unindo estatísticos, cientistas da computação e especialistas do domínio do negócio \cite{morettin2019introducao}.

\section{Avaliação de Interface e Desempenho do Sistema}
\label{sec:Teste}

Para avaliar a efetividade de um sistema \textit{desktop}, é de suma importância a realização de testes focados na estabilidade e no desempenho da aplicação. A validação consiste em analisar a melhoria contínua do Painel Parkinson, observando fatores técnicos como o tempo de carregamento de dados, a redução de latência e a eficiência na apresentação visual das informações, garantindo que o sistema responda rapidamente aos comandos.

Para uma efetiva validação, é necessário submeter a ferramenta \textit{desktop} a simulações que reflitam o volume de dados de um cenário de uso real. Ao longo dessa execução, a coleta de métricas de desempenho garante a comprovação de que a arquitetura suporta a exigência do ambiente clínico.

Segundo Hewett et al. (1992), a área de Interação Humano-Computador (IHC) está relacionada ao projeto, implementação e avaliação de sistemas computacionais interativos voltados ao uso humano. Nesse contexto, aspectos como contexto de uso, características dos usuários e arquitetura da interface influenciam diretamente na experiência de utilização do sistema \cite{ihc2022}.

Além disso, os estudos relacionados à Experiência do Usuário (\textit{User Experience} -- UX) destacam que fatores como facilidade de uso, eficiência, acessibilidade e organização visual influenciam diretamente na aceitação e utilização de sistemas computacionais em ambientes profissionais. Dessa forma, a adaptação do Painel Parkinson para o ambiente \textit{desktop} busca não apenas uma mudança estrutural da aplicação, mas também melhorias de desempenho que impactam positivamente a experiência de uso dos profissionais da saúde.

Segundo definem Few e Edge (2007 apud \citeonline[p. 20]{oliveira2023}), "Um painel é uma exibição visual das informações mais importantes necessárias para atingir um ou mais objetivos; consolidado e organizado em uma única tela para que as informações possam ser monitoradas rapidamente." Nesse sentido, o painel informativo deve ser submetido a rigorosos testes de validação técnica para garantir que a robustez do sistema permita que os profissionais envolvidos no acompanhamento do paciente consigam abstrair os conteúdos propostos de forma fluida e sem interrupções.

\section{Protocolo de Mapeamento Sistemático da Literatura}
\label{sec:ProtocoloSistematico}

A fim de identificar o estado da arte e as lacunas tecnológicas acerca do monitoramento clínico da Doença de Parkinson (DP), realizou-se um Mapeamento Sistemático da Literatura. Conforme apontado por \citeonline{demuth2024scoping}, essa metodologia é fundamental para mapear abordagens algorítmicas e verificar a maturidade clínica de sistemas de suporte à decisão. A progressão variável da DP impõe desafios significativos ao manejo clínico tradicional, baseado em consultas presenciais esporádicas, tornando necessária a implementação de estratégias de monitoramento contínuo \citeonline{mahboobeh2021dynamic}.

\section{Evolução das Soluções e Base Estrutural}
\label{sec:EvolucaoSolucoes}

O mapeamento revelou uma transição de modelos teóricos para implementações digitais disruptivas. Iniciativas como o projeto SIDABI, explorado por \citeonline{folador2021}, foram pioneiras na centralização de dados biomédicos e apoio a grupos de pesquisa. A partir dessa base, surgiram soluções focadas na visualização clínica, exemplificadas pelo Painel Parkinson, cujos estudos de \citeonline{oliveira2023} e \citeonline{rabelo2025} expandiram o uso de \textit{dashboards} interativos e mecanismos de suporte à decisão baseados em estatísticas temporais \citeonline{santos2023desenvolvimento}.

A implementação tecnológica dessas soluções fundamenta-se em ecossistemas modulares que integram aplicações voltadas ao registro de testes de sintomas e controle medicamentoso a dashboards centralizados para o uso médico \citeonline{mahboobeh2021dynamic}. Arquiteturas modernas, como o padrão Model-View-Controller (MVC) e serviços de \textit{Backend} via APIs REST, garantem a segurança, escalabilidade e o processamento em tempo real dos dados \citeonline{mahboobeh2021dynamic}.

\section{Lacunas Identificadas e Proposta}
\label{sec:Lacunas}

Apesar do avanço, o mapeamento sistemático evidenciou limitações críticas: a dependência de navegadores e a necessidade de conexão constante com a internet restringem a praticidade no ambiente clínico. Além disso, a literatura aponta que a eficácia diagnóstica é potencializada por funcionalidades avançadas de "contextualização", onde métricas individuais são comparadas a populações de referência, e pelo uso de sensores vestíveis para a coleta quantitativa de sintomas motores, como a rigidez de pulso \citeonline{demuth2024scoping, mahboobeh2021dynamic}. Diante disso, a proposta deste trabalho busca preencher essa lacuna por meio da adaptação do Painel Parkinson para uma arquitetura \textit{desktop}, visando elevar o desempenho, a acessibilidade e a praticidade da ferramenta no cotidiano hospitalar \citeonline{santos2023desenvolvimento}.

\subsection{Sistemas Similares e Trabalhos Correlatos}
\label{subsec:sistemas_similares_correlatos}

Com base no mapeamento sistemático da literatura e na análise crítica das soluções desenvolvidas no âmbito do Grupo de Pesquisa em Computação Aplicada (GPCA), foi possível mapear a evolução das ferramentas tecnológicas de suporte ao acompanhamento da Doença de Parkinson (DP). Embora as soluções anteriores tenham estabelecido as bases conceituais para a gestão de dados clínicos, a análise de suas respectivas implementações revela um conjunto de lacunas arquiteturais e falhas de usabilidade que comprometem sua aplicação prática contínua.

No cenário de soluções correlatas globais, destacam-se iniciativas consolidadas voltadas ao monitoramento de sintomas. O estudo \textit{Fox Insight} (2017) baseia-se em avaliações longitudinais \textit{online} por meio de formulários digitais autorrelatados pelos pacientes, porém carece de integração direta com bancos de dados biomédicos de referência local. Por sua vez, a plataforma comercial \textit{Empatica PKG} (2025) utiliza sensores vestíveis (\textit{wearables}) de nível médico para o registro contínuo de tremores e sintomas motores, processando relatórios em infraestrutura de nuvem proprietária. Apesar da robustez de ambas as soluções, observa-se que sua operação é estritamente subordinada ao modelo de computação em nuvem, o que gera custos recorrentes de processamento, alta dependência de conectividade ininterrupta e restrições de compatibilidade com ecossistemas acadêmicos abertos.

No contexto das pesquisas desenvolvidas internamente pelo GPCA, a evolução tecnológica deu-se a partir do projeto SIDABI (Sistema Integrado de Dados Biomédicos) \cite{folador2021}, que estruturou o armazenamento, gerenciamento e segurança de dados clínicos brutos relacionados à DP. Com base nessa infraestrutura, \cite{oliveira2021} desenvolveram o primeiro protótipo de \textit{dashboard} em ambiente \textit{web} focado na visualização gráfica das escalas clínicas da \textit{Unified Parkinson's Disease Rating Scale} (MDS-UPDRS). Posteriormente, o sistema de apresentação de dados proposto por \cite{costa2022sistema} integrou a visualização das sessões de coletas feitas por sensores inerciais de acelerometria e giroscopia. Mais recentemente, \cite{oliveira2023} expandiram a plataforma por meio do desenvolvimento de painéis informativos para avaliar a influência de terapias complementares no estado emocional dos pacientes, utilizando escalas clínicas integradas.

Apesar da contribuição científica inegável das versões \textit{web} anteriores, a consolidação desse ecossistema em ambiente clínico revelou gargalos operacionais críticos que inviabilizam a sua adoção fluida no cotidiano hospitalar. O primeiro gargalo refere-se à latência de processamento e renderização de dados biométricos pesados. Como evidenciado por \cite{costa2022sistema}, o processamento de arquivos volumosos de sensores em nuvem (hospedados em instâncias de baixa escala como AWS EC2 e AWS RDS) degrada o desempenho da aplicação, gerando travamentos no navegador do usuário e comprometendo a visualização de gráficos dinâmicos de alta frequência devido à latência de rede entre o cliente e o servidor.

Do ponto de vista da infraestrutura e viabilidade econômica, o funcionamento ininterrupto da aplicação em nuvem sob demanda eleva substancialmente os custos operacionais com hospedagem de servidores e manutenção de bancos de dados gerenciados, como o PostgreSQL Aurora, conforme pontuado por \cite{oliveira2023}. Ademais, a arquitetura estritamente baseada na \textit{web} impõe uma dependência crítica de conectividade estável com a internet. Em ambientes clínicos e hospitalares, sujeitos a instabilidades severas de rede, essa dependência pode interromper os atendimentos médicos e impossibilitar o acesso em tempo real aos históricos dos pacientes, o que representa um fator de risco à continuidade assistencial.

Essas limitações arquiteturais refletiram-se diretamente em falhas de usabilidade severas identificadas pelos usuários finais. No estudo clínico de usabilidade conduzido por \cite{rabelo2025} junto a 10 profissionais da saúde (fisioterapeutas e profissionais de educação física), foram mapeadas falhas sistêmicas recorrentes de navegação e exibição no painel \textit{web}:

\begin{itemize}
    \item \textbf{Perda de estado de seleção (IDP1, IDP2, IDP5):} A aplicação não conseguia manter parâmetros de filtragem definidos pelo profissional (tais como a condição ``com estímulo'' ou ``sem estímulo'') após a atualização de páginas ou retorno a telas anteriores, exigindo reconfigurações constantes;
    \item \textbf{Renderização incorreta e perda de pontos gráficos (IDP6):} Certas rotinas de \textit{zoom} em gráficos dinâmicos de tremor faziam com que pontos de dados sumissem da tela, comprometendo a precisão da leitura médica;
    \item \textbf{Carregamento incorreto de dados e indução ao erro (IDP8):} O sistema pré-carregava informações antigas ou inconsistentes na interface UPDRS antes da validação completa da seleção do profissional de saúde, induzindo o especialista a interpretar dados equivocados de outros pacientes;
    \item \textbf{Navegação redundante e duplicação de histórico (IDP4):} O uso do botão de voltar nativo do navegador gerava redirecionamentos cíclicos, forçando recarregamentos desnecessários e duplicando registros no histórico de navegação, elevando o tempo de conclusão das tarefas clínicas.
\end{itemize}

Frente a esses desafios, a transição arquitetural do Painel Parkinson para uma estrutura local multiplataforma via \textit{desktop} surge como a solução tecnológica definitiva para mitigar tais gargalos. Ao envelopar o \textit{front-end} em VueJS utilizando o \textit{framework} Electron.js, e integrando localmente o \textit{back-end} em Python (Flask) com o banco de dados PostgreSQL, os seguintes benefícios diretos são proporcionados ao cenário clínico:

\begin{itemize}
    \item \textbf{Latência Zero e Renderização Gráfica Avançada:} O processamento local elimina os tempos de tráfego de dados biométricos de alta amostragem (sensores inerciais e matrizes densas), garantindo renderização visual instantânea e responsividade de gráficos clínicos de tremor, sem sobrecarregar a memória RAM ou congelar a interface do usuário;
    \item \textbf{Independência de Conexão (Suporte Offline):} A persistência de dados local permite o pleno funcionamento do sistema em modo completamente \textit{offline}, assegurando que o profissional de saúde realize consultas e coletas ininterruptas em laboratórios e enfermarias, independentemente de instabilidades de rede;
    \item \textbf{Otimização de Custos e Escalabilidade Acadêmica:} A descentralização da infraestrutura computacional elimina a necessidade de manutenção de instâncias de processamento robustas em servidores de nuvem, reduzindo os custos recorrentes de hospedagem e viabilizando a distribuição sustentável do \textit{software} para múltiplos centros de pesquisa parceiros;
    \item \textbf{Conformidade com a LGPD e Privacidade de Dados:} Em estrita consonância com a Lei Geral de Proteção de Dados (LGPD), a arquitetura local permite que os dados sensíveis dos pacientes sejam armazenados e processados sob o controle direto do hospital ou laboratório de coleta. A sincronização com a base central de pesquisa do SIDABI é realizada de forma automática em segundo plano (\textit{background}) por meio de protocolos de anonimização (em conformidade com o Art. 13 da Lei nº 13.709/2018), eliminando o risco de exposição de informações em trânsito e mitigando vulnerabilidades cibernéticas.
\end{itemize}

Dessa forma, a presente solução proposta desvia o foco de uma abordagem científica abstrata e apresenta um benefício tecnológico palpável, convertendo o Painel Parkinson em um Sistema de Apoio à Decisão Clínica (SADC) verdadeiramente robusto, rápido, financeiramente viável e seguro para a prática médica cotidiana.

\begin{table}[htbp]
\centering
\caption{Comparativo entre sistemas similares e a proposta atual.}
\label{tab:sistemas_similares}
\begin{tabular}{|p{3.8cm}|p{5.2cm}|p{6.5cm}|}
\hline
\textbf{Sistema} & \textbf{Objetivo} & \textbf{Limitações / Oportunidades} \\ \hline
\textbf{Fox Insight (2017)} & Estudo longitudinal \textit{online} para coleta de dados de rotina por meio de questionários virtuais. & Baseado em dados autorrelatados via \textit{web}, sem integração direta com banco de dados biomédicos locais. \\ \hline
\textbf{SIDABI (2021)} & Gerenciamento, centralização e segurança de dados clínicos brutos da Doença de Parkinson. & Foco estrito em armazenamento de dados estruturados e de coletas, sem camada avançada de visualização ou suporte à decisão. \\ \hline
\textbf{Dashboard Web (2021)} & Visualização gráfica inicial para interpretação de escalas clínicas (MDS-UPDRS) do banco SIDABI. & Dependência exclusiva de navegador \textit{web}, dependência de internet e instâncias de nuvem limitadas. \\ \hline
\textbf{Visualização de Sensores (2022)} \newline (Lucas Costa) & Apresentação gráfica de dados coletados por sensores inerciais (acelerometria e giroscopia). & Elevada latência de processamento de arquivos CSV volumosos na nuvem e travamentos de renderização dinâmica no navegador. \\ \hline
\textbf{Painel Parkinson (2023)} \newline (Lucas e Samuel) & Visualização de dados emocionais e acompanhamento de terapias complementares da DP. & Arquitetura \textit{web} sujeita a instabilidades de rede, altos custos de nuvem recorrentes e perda de estado de filtros clínicos. \\ \hline
\textbf{Empatica PKG (2025)} & Monitoramento remoto e contínuo de sintomas motores via \textit{wearables} de nível médico. & Solução comercial proprietária de alto custo, com dependência de conectividade constante à nuvem privada e sem integração aberta com bases locais. \\ \hline
\textbf{Painel Parkinson Desktop (2026) [Proposta Atual]} & \textbf{Apoio à decisão clínica local com processamento de alto desempenho e gráficos instantâneos.} & \textbf{Alta performance local (latência zero), operação 100\% offline para consultórios, custo zero de infraestrutura em nuvem e conformidade nativa com a LGPD.} \\ \hline
\end{tabular}
\vspace{0.2cm}

\small{\textit{Fonte: Elaborado pelos autores com base em Michael J. Fox Foundation (2017), Folador (2021), Oliveira e Conceição (2021), Costa (2022), Oliveira e Gouvea (2023) e Empatica (2025).}}
\end{table}
